\documentclass[11pt,a4paper]{article}

% --- Encoding ---
\usepackage[utf8]{inputenc}
\usepackage[T1]{fontenc}
\usepackage[spanish,es-noshorthands]{babel}
\usepackage{lmodern}

% --- Layout ---
\usepackage[margin=2.5cm]{geometry}
\usepackage{parskip}
\setlength{\parindent}{0pt}
\setlength{\parskip}{0.5em plus 0.2em minus 0.1em}

% --- Tables and graphics ---
\usepackage{booktabs}
\usepackage{array}
\usepackage{enumitem}
\usepackage{xcolor}
\usepackage{graphicx}
\usepackage{tikz}
\usetikzlibrary{positioning,shapes.geometric,arrows.meta,fit,backgrounds,calc,decorations.pathreplacing}
\usepackage{caption}
\usepackage{subcaption}
\usepackage{amsmath,amssymb}

% --- Soporte para caracteres Unicode usados en formulas en texto plano ---
\usepackage{newunicodechar}
\newunicodechar{→}{\ensuremath{\rightarrow}}
\newunicodechar{≥}{\ensuremath{\geq}}
\newunicodechar{≤}{\ensuremath{\leq}}
\newunicodechar{×}{\ensuremath{\times}}
\newunicodechar{÷}{\ensuremath{\div}}
\newunicodechar{≈}{\ensuremath{\approx}}
\newunicodechar{⌈}{\ensuremath{\lceil}}
\newunicodechar{⌉}{\ensuremath{\rceil}}
\newunicodechar{✓}{\ensuremath{\checkmark}}

% --- Code listings ---
\usepackage{listings}
\definecolor{codegray}{rgb}{0.96,0.96,0.96}
\definecolor{codeborder}{rgb}{0.85,0.85,0.85}
\definecolor{codecomment}{rgb}{0.30,0.55,0.30}
\definecolor{codekeyword}{rgb}{0.10,0.30,0.70}
\definecolor{codestring}{rgb}{0.65,0.15,0.10}

\lstdefinestyle{cstyle}{
    language=C,
    backgroundcolor=\color{codegray},
    basicstyle=\ttfamily\footnotesize,
    breaklines=true,
    frame=single,
    rulecolor=\color{codeborder},
    showstringspaces=false,
    keywordstyle=\color{codekeyword}\bfseries,
    commentstyle=\color{codecomment}\itshape,
    stringstyle=\color{codestring},
    aboveskip=1em,
    belowskip=1em,
}

\lstdefinestyle{shellstyle}{
    language=bash,
    backgroundcolor=\color{codegray},
    basicstyle=\ttfamily\small,
    breaklines=true,
    frame=single,
    rulecolor=\color{codeborder},
    showstringspaces=false,
    keywordstyle=\color{codekeyword}\bfseries,
    commentstyle=\color{codecomment}\itshape,
    stringstyle=\color{codestring},
    aboveskip=1em,
    belowskip=1em,
    columns=fullflexible,
    upquote=true,
}

\lstdefinestyle{textstyle}{
    backgroundcolor=\color{codegray},
    basicstyle=\ttfamily\footnotesize,
    breaklines=true,
    frame=single,
    rulecolor=\color{codeborder},
    showstringspaces=false,
    aboveskip=1em,
    belowskip=1em,
    columns=fullflexible,
}

\lstset{style=textstyle}

% --- Definition / Idea / Important boxes ---
\usepackage[most]{tcolorbox}
\newtcolorbox{idea}[1][]{
    colback=orange!5!white,
    colframe=orange!60!black,
    fonttitle=\bfseries,
    title=Idea intuitiva,
    sharp corners=southwest,
    rounded corners=northwest,
    arc=2pt,
    boxrule=0.6pt,
    left=8pt,right=8pt,top=4pt,bottom=4pt,
    #1
}
\newtcolorbox{importante}[1][]{
    colback=red!5!white,
    colframe=red!50!black,
    fonttitle=\bfseries,
    title=Importante,
    sharp corners=southwest,
    rounded corners=northwest,
    arc=2pt,
    boxrule=0.6pt,
    left=8pt,right=8pt,top=4pt,bottom=4pt,
    #1
}
\newtcolorbox{respuesta}[1][]{
    colback=green!5!white,
    colframe=green!40!black,
    fonttitle=\bfseries,
    title=Respuesta,
    sharp corners=southwest,
    rounded corners=northwest,
    arc=2pt,
    boxrule=0.6pt,
    left=8pt,right=8pt,top=4pt,bottom=4pt,
    #1
}

% --- Hyperlinks ---
\usepackage[colorlinks=true,
            linkcolor=blue!60!black,
            urlcolor=blue!50!black,
            citecolor=blue!60!black]{hyperref}

% --- Color palette ---
\definecolor{c1}{HTML}{4C72B0}
\definecolor{c2}{HTML}{DD8452}
\definecolor{c3}{HTML}{55A868}
\definecolor{c4}{HTML}{8172B2}
\definecolor{cred}{HTML}{C44E52}
\definecolor{cfree}{HTML}{8C8C8C}

% --- Title block ---
\title{
    \vspace{-2.5em}
    \textbf{Examen Parcial 1 --- 2025} \\[0.3em]
    \large Sistemas Operativos
}
\author{
    Tomás Rodeghiero \\
    \small UNRC --- Universidad Nacional de Río Cuarto
}
\date{2025}

\begin{document}

\maketitle

\begin{abstract}
\noindent
Resolución completa del Parcial 1 del año 2025 de la cátedra de
Sistemas Operativos (UNRC). El parcial consta de siete ejercicios
que cubren: (i) IPC síncrono en un microkernel, contando traps y
context switches y comparándolos con un kernel monolítico; (ii)
preguntas sobre pipes y syscalls que crean procesos; (iii)
implementación de \texttt{lock}/\texttt{unlock} sin instrucciones
atómicas y análisis de fallas en monoprocesador y multiprocesador;
(iv) verdadero/falso sobre \emph{multilevel feedback queue} (MLFQ);
(v) ejemplo de deadlock con \texttt{lock}/\texttt{unlock}; (vi)
verdadero/falso sobre paginación; (vii) cálculos sobre una
arquitectura de 32 bits con páginas de 4\,MB (cantidad de páginas,
interpretación de la dirección virtual, tabla de un proceso y
detección de page faults). Cada ejercicio se desarrolla con la
estructura \textbf{Enunciado $\to$ Análisis $\to$ Desarrollo $\to$
Respuesta}.
\end{abstract}

\tableofcontents
\bigskip

% =====================================================================
\section{Ejercicio 1 --- IPC síncrono en microkernel (1 pto)}
% =====================================================================

\textbf{Enunciado.} En un SO con diseño microkernel, un proceso $U$
realiza una llamada al sistema de la forma
\texttt{send(process\_manager, "getpid")}. El microkernel reenvía el
mensaje al \texttt{process\_manager}. El mecanismo de mensajes es
\textbf{síncrono} y se basa en interrupciones por software, de modo
que cada operación \texttt{send} o \texttt{receive} causa un trap al
microkernel.

\textbf{(a)} Determinar la cantidad de traps y de cambios de
contexto entre procesos que requiere esta llamada al sistema hasta
que el proceso recibe la respuesta.

\textbf{(b)} Compararlo con un kernel monolítico.

\subsection*{Análisis: arquitectura microkernel}

En un microkernel los \emph{servicios del SO} (gestor de memoria,
file system, drivers, gestor de procesos, etc.) corren como procesos
en \textbf{espacio de usuario}. La única función del microkernel
mínimo es \emph{enviar mensajes} entre procesos y gestionar la
planificación y trampas de hardware.

Para que $U$ obtenga su PID, no puede preguntar directamente al
kernel (como en monolítico): tiene que \emph{enviar un mensaje} al
proceso servidor llamado \texttt{process\_manager} (PM), que vive en
modo usuario y responde.

\subsection*{1.a --- Conteo de traps y context switches}

Trazamos el camino completo de la llamada:

\begin{center}
\begin{tikzpicture}[
    every node/.style={font=\ttfamily\small,align=center},
    box/.style={draw,rounded corners,minimum width=1.6cm,minimum height=0.8cm},
    user/.style={box,fill=c1!25},
    kern/.style={box,fill=c2!25},
    pm/.style={box,fill=c3!25},
    arr/.style={->,thick}
]
\node[user] (u1) at (0,3)  {$U$};
\node[kern] (k1) at (3,3)  {$\mu$kernel};
\node[pm]   (p1) at (6,3)  {PM};
\node[kern] (k2) at (3,0.5){$\mu$kernel};
\node[user] (u2) at (0,-2) {$U$};

\draw[arr] (u1.east) -- node[above,font=\scriptsize]{trap 1 (\texttt{send})} (k1.west);
\draw[arr,dashed] (k1.east) -- node[above,font=\scriptsize]{ctx switch 1} (p1.west);
\draw[arr] (p1.south) -- node[right,font=\scriptsize]{trap 2 (\texttt{send} reply)} (k2.north east);
\draw[arr,dashed] (k2.west) -- node[above,font=\scriptsize]{ctx switch 2} ++(-3,0) |- (u2.east);
\end{tikzpicture}
\end{center}

\textbf{Secuencia mínima} (send sincrónico, estilo Mach/L4/QNX,
\emph{send-and-receive-reply}):

\begin{enumerate}[leftmargin=*,nosep]
    \item $U$ ejecuta \texttt{send(PM, "getpid")} en modo usuario.
        \textbf{Trap 1} al microkernel.
    \item El microkernel encola el mensaje para PM. $U$ queda
        bloqueado esperando la respuesta.
    \item PM estaba previamente bloqueado en un \texttt{receive}
        (esperando peticiones). El microkernel lo despierta y
        ejecuta un \textbf{cambio de contexto 1} ($U \to$ PM).
    \item PM en modo usuario procesa la petición (computa el PID)
        y ejecuta \texttt{send(U, "$pid$")} para devolver la
        respuesta. \textbf{Trap 2} al microkernel.
    \item El microkernel entrega la respuesta a $U$ (que estaba
        bloqueado). \textbf{Cambio de contexto 2} (PM $\to U$).
    \item $U$ reanuda y obtiene la respuesta.
\end{enumerate}

\begin{respuesta}
\textbf{(a)} La llamada requiere:

\begin{itemize}[leftmargin=*,nosep]
    \item \textbf{2 traps} al microkernel (uno por cada \texttt{send}).
    \item \textbf{2 cambios de contexto} entre procesos
        ($U \to$ PM y PM $\to U$).
\end{itemize}

\textbf{Nota}: si la API separara \texttt{send} y \texttt{receive}
en syscalls independientes (estilo POSIX message queues), el conteo
sería de 3 o 4 traps. La cantidad de cambios de contexto se mantiene
en 2.
\end{respuesta}

\subsection*{1.b --- Comparación con kernel monolítico}

En un kernel monolítico, \texttt{getpid()} es una syscall provista
directamente por el kernel:

\begin{enumerate}[leftmargin=*,nosep]
    \item $U$ ejecuta \texttt{getpid()}. \textbf{Trap 1} al kernel.
    \item El kernel lee el PID del PCB del proceso actual y retorna.
    \item $U$ reanuda con el valor.
\end{enumerate}

\textbf{Total}: \textbf{1 trap, 0 cambios de contexto}.

\begin{idea}
La comparación muestra el \emph{overhead inherente} del modelo
microkernel para syscalls simples: para obtener un valor que en
monolítico se resuelve con $\sim 50$\,ns de trap, el microkernel
agrega $2 \times$ cambio de contexto $+ 2 \times$ trap, fácilmente
$1$\,$\mu$s o más. Por eso los microkernels (Mach, MINIX) son
\emph{conceptualmente elegantes} pero \emph{históricamente más
lentos} que los monolíticos para cargas típicas. L4 y derivados
modernos (seL4, Fiasco) han optimizado mucho este camino, pero la
diferencia estructural sigue existiendo.
\end{idea}

% =====================================================================
\section{Ejercicio 2 --- Comandos y syscalls (1 pto)}
% =====================================================================

\subsection*{2.a --- ``\texttt{p1 | p2} hace que p2 se bloquee hasta que p1 escriba''}

\begin{respuesta}
\textbf{V}. El operador \texttt{|} conecta el \texttt{stdout} de
\texttt{p1} con el \texttt{stdin} de \texttt{p2} mediante un pipe del
kernel. Si \texttt{p2} llama a \texttt{read(stdin, \ldots)} y el pipe
está vacío, el \texttt{read} bloquea al proceso (estado
\texttt{SLEEPING}) hasta que haya datos disponibles. Cuando \texttt{p1}
escriba, el kernel despierta a \texttt{p2} y el \texttt{read} retorna.
\end{respuesta}

\subsection*{2.b --- ``Si p1 se planifica primero, p1 se bloqueará al escribir''}

\begin{respuesta}
\textbf{F}. \texttt{p1} \emph{no} se bloquea al escribir mientras
haya espacio en el buffer interno del pipe (\texttt{PIPE\_BUF}, en
Linux típicamente $64$\,KB). El escritor solo se bloquea cuando el
buffer está \emph{lleno} y nadie lo está consumiendo. Si \texttt{p1}
escribe poco (o si \texttt{p2} lee a un ritmo razonable), no hay
bloqueo del escritor: la idea del pipe es justamente desacoplar
productor y consumidor.

\textbf{Excepción}: si \texttt{p1} es CPU-bound y produce gigas de
datos sin que \texttt{p2} consuma, sí terminará bloqueándose.
\end{respuesta}

\subsection*{2.c --- Eliminar duplicados y ordenar sin temporales}

Comando idiomático:

\begin{lstlisting}[style=shellstyle]
sort archivo | uniq > email-final.txt
\end{lstlisting}

\quad O en una sola pasada con \texttt{sort -u}:

\begin{lstlisting}[style=shellstyle]
sort -u archivo > email-final.txt
\end{lstlisting}

\begin{respuesta}
La forma canónica es \verb|sort archivo | uniq > email-final.txt|.

\textbf{Por qué funciona sin archivos auxiliares:}

\begin{itemize}[leftmargin=*,nosep]
    \item \texttt{sort} lee el archivo, lo ordena y vuelca a
        \texttt{stdout}.
    \item El pipe conecta esa salida con \texttt{uniq} como entrada.
    \item \texttt{uniq} requiere su entrada \emph{ordenada} para
        funcionar (elimina duplicados \emph{adyacentes}); por eso
        viene después de \texttt{sort}.
    \item \texttt{>} redirige la salida final al archivo
        \texttt{email-final.txt}.
\end{itemize}

Los datos viajan por buffers en RAM del kernel; no hay archivos
intermedios en disco.
\end{respuesta}

\subsection*{2.d --- ``\texttt{fork} y \texttt{exec} son las únicas que crean procesos''}

\begin{respuesta}
\textbf{F}. \texttt{exec} \emph{no crea un proceso nuevo}: reemplaza
la imagen del proceso actual con la del nuevo programa (mismo PID,
mismo descriptor). La única syscall que efectivamente \emph{crea} un
nuevo proceso es \texttt{fork()} (y sus variantes
\texttt{vfork()}/\texttt{clone()} en Linux, internamente las tres
descansan sobre el mismo mecanismo).

El patrón clásico \textbf{fork+exec} usa las dos juntas: \texttt{fork}
crea el hijo, y \texttt{exec} lo transforma en otro programa. Pero la
\emph{creación} es solo del \texttt{fork}.
\end{respuesta}

% =====================================================================
\section{Ejercicio 3 --- \texttt{lock}/\texttt{unlock} sin atomicidad (2 ptos)}
% =====================================================================

\textbf{Enunciado.} Dar pseudo-código de \texttt{lock(\&lk)} y
\texttt{unlock(\&lk)} que \emph{no} deshabilitan ni habilitan
interrupciones, asumiendo un sistema mono-procesador.

\subsection*{Pseudo-código (incorrecto a propósito)}

\begin{lstlisting}[style=cstyle]
typedef int lock_t;     // 0 = libre, 1 = tomado

void lock(lock_t *lk) {
    while (*lk == 1)     // 1) test
        ;                //    busy-wait
    *lk = 1;             // 2) set
}

void unlock(lock_t *lk) {
    *lk = 0;
}
\end{lstlisting}

\textbf{El problema fundamental}: la pareja ``test luego set''
(líneas 1 y 2 de \texttt{lock}) \textbf{no es atómica}: entre la
lectura de \texttt{*lk} y la escritura de \texttt{*lk = 1} puede
ocurrir una interrupción que dé control a otro thread, el cual entre
también a la sección crítica.

\subsection*{3.a --- Traza que viola la exclusión mutua (monoprocesador)}

Sean dos threads $T_1$ y $T_2$ en un único procesador. Inicialmente
\texttt{lk = 0}.

\begin{center}
\small
\begin{tabular}{@{}clp{6cm}p{3.5cm}@{}}
\toprule
$t$ & Tarea ejecutando & Operación & Estado \\
\midrule
$t_1$ & $T_1$ & \texttt{while (*lk == 1)} $\to$ lee $0$, sale & \texttt{lk} = 0 \\
$t_2$ & \multicolumn{2}{l}{\textbf{interrupción del timer}: scheduler cambia a $T_2$} & \\
$t_3$ & $T_2$ & \texttt{while (*lk == 1)} $\to$ lee $0$, sale & \texttt{lk} = 0 \\
$t_4$ & $T_2$ & \texttt{*lk = 1} & \texttt{lk} = 1 \\
$t_5$ & $T_2$ & \textbf{ENTRA a la sección crítica} & \\
$t_6$ & \multicolumn{2}{l}{\textbf{interrupción del timer}: scheduler cambia a $T_1$} & \\
$t_7$ & $T_1$ & \texttt{*lk = 1} (siguiendo desde donde quedó) & \texttt{lk} = 1 \\
$t_8$ & $T_1$ & \textbf{ENTRA a la sección crítica} & \\
\bottomrule
\end{tabular}
\end{center}

\textbf{Resultado}: ambos threads están simultáneamente dentro de la
sección crítica. \textbf{No se garantiza exclusión mutua.}

La raíz del problema: $T_1$ \emph{ya había decidido} entrar
(\texttt{while} salió con \texttt{lk = 0}) pero \emph{aún no había
escrito} \texttt{lk = 1}. La interrupción dejó el ``lock'' en un
estado inconsistente: $T_1$ comprometido pero aún no marcado.

\subsection*{3.b --- Multiprocesador, aún con IRQs deshabilitadas}

Ahora supongamos que \texttt{lock} \emph{sí} deshabilita IRQs en la
CPU local antes del test-set, y las habilita después. Pero el sistema
es multiprocesador (varias CPUs corriendo en paralelo).

\textbf{Problema}: deshabilitar IRQs en la CPU $A$ no impide que la
CPU $B$ ejecute simultáneamente. Las dos CPUs pueden estar en la
sección de test-set al mismo tiempo, y como ``test luego set'' no es
atómico a nivel de memoria, ambas terminan entrando.

\textbf{Traza:}

\begin{center}
\small
\begin{tabular}{@{}clll@{}}
\toprule
$t$ & CPU $A$ (corre $T_1$) & CPU $B$ (corre $T_2$) & \texttt{lk} \\
\midrule
$t_1$ & \texttt{cli} (deshabilita IRQs en $A$) & \texttt{cli} (deshabilita IRQs en $B$) & 0 \\
$t_2$ & lee \texttt{*lk} $\to$ 0, sale del while & lee \texttt{*lk} $\to$ 0, sale del while & 0 \\
$t_3$ & \texttt{*lk = 1} & \texttt{*lk = 1} (lo mismo, ya estaba 1) & 1 \\
$t_4$ & \textbf{entra SC} & \textbf{entra SC}  & 1 \\
\bottomrule
\end{tabular}
\end{center}

Ambas CPUs leen \texttt{lk = 0} casi simultáneamente; ambas escriben
\texttt{*lk = 1}; ambas entran a la sección crítica. \textbf{Falla.}

\begin{importante}
\textbf{Conclusión}: deshabilitar interrupciones funciona como
mecanismo de exclusión \emph{solo en monoprocesador}, porque allí la
preemption es la única fuente de concurrencia. En multiprocesador,
\textbf{además} hay paralelismo físico entre CPUs, y deshabilitar
IRQs no afecta a otras CPUs.

La solución para multiprocesador es usar \textbf{instrucciones
atómicas del hardware}:

\begin{itemize}[leftmargin=*,nosep]
    \item \textbf{Test-And-Set} (TAS): \texttt{xchg} en x86;
        actualiza el valor y retorna el anterior en una sola
        operación que el bus de memoria garantiza atómica.
    \item \textbf{Compare-And-Swap} (CAS): \texttt{cmpxchg};
        condicional, base de muchos algoritmos lock-free.
    \item \textbf{LL/SC}: \texttt{ldrex/strex} (ARM),
        \texttt{lr/sc} (RISC-V); par de instrucciones cuya
        atomicidad la detecta el cache coherence.
\end{itemize}

Con cualquiera de estas, el ``test luego set'' se vuelve atómico y
la exclusión mutua queda garantizada.
\end{importante}

% =====================================================================
\section{Ejercicio 4 --- V/F sobre MLFQ (1.5 ptos)}
% =====================================================================

\textbf{Enunciado.} Un SO usa MLFQ con la siguiente regla: si un
proceso \emph{no} usó todo su quantum, se le \textbf{sube} la
prioridad (al nivel superior). Si lo usó completo, se le \textbf{baja}
la prioridad. ¿V/F?

\subsection*{Análisis: qué procesos suben y cuáles bajan}

\begin{itemize}[leftmargin=*,nosep]
    \item Un proceso \textbf{CPU-bound} ejecuta su quantum entero
        sin bloquearse (no hace I/O). Por la regla, \textbf{baja} de
        prioridad.
    \item Un proceso \textbf{I/O-bound} típicamente se bloquea en una
        operación de I/O (\texttt{read}, \texttt{recv}, etc.) antes
        de agotar el quantum. Por la regla, \textbf{sube} de prioridad.
\end{itemize}

Resultado: los procesos interactivos (I/O-bound) terminan en los
niveles altos del MLFQ y se planifican \emph{antes}; los CPU-bound
quedan relegados a los niveles bajos.

\subsection*{4.a --- ``El scheduler prioriza a los procesos CPU-bound''}

\begin{respuesta}
\textbf{F}. Es exactamente al revés. Los procesos CPU-bound
consumen siempre todo su quantum y por eso \emph{bajan} de prioridad;
quedan en los niveles más bajos del MLFQ y se planifican menos.
\end{respuesta}

\subsection*{4.b --- ``El scheduler prioriza a los procesos I/O-bound''}

\begin{respuesta}
\textbf{V}. Los I/O-bound se bloquean antes de gastar su quantum,
así que por la regla suben de prioridad. Tienden a quedar en los
niveles altos, lo que les da menor tiempo de respuesta --- ideal para
procesos interactivos.
\end{respuesta}

\subsection*{4.c --- ``El objetivo es maximizar el throughput''}

\begin{respuesta}
\textbf{F}. El objetivo principal de MLFQ no es maximizar
\emph{throughput} (cantidad de procesos completados por unidad de
tiempo). Para eso conviene FCFS o SJF (que minimizan el promedio de
turnaround sin la sobrecarga de muchos cambios de contexto).

MLFQ se diseñó para optimizar el \textbf{tiempo de respuesta} de
los procesos interactivos manteniendo \emph{fairness}: la prioridad
adaptativa hace que los procesos I/O-bound (que esperan input del
usuario o de la red) se ejecuten apenas tienen trabajo, sin sacrificar
totalmente a los CPU-bound.
\end{respuesta}

\subsection*{4.d --- ``Minimizar el tiempo medio de espera de los I/O-bound''}

\begin{respuesta}
\textbf{V}. Es exactamente el objetivo: que los procesos
orientados a I/O (interactivos) reciban CPU rápidamente cuando
están listos para ejecutar. Esto reduce su tiempo de espera en la
cola \texttt{READY} y, por ende, mejora la latencia percibida por el
usuario.

Un usuario tipeando en un editor o esperando una respuesta del
servidor web nota fundamentalmente el tiempo entre que termina su
input y la pantalla responde. MLFQ optimiza precisamente esa métrica.
\end{respuesta}

% =====================================================================
\section{Ejercicio 5 --- Ejemplo de deadlock con \texttt{lock}/\texttt{unlock} (1.5 ptos)}
% =====================================================================

\textbf{Enunciado.} Dar un ejemplo de dos tareas (procesos o threads)
que produzcan un deadlock usando las primitivas \texttt{lock} y
\texttt{unlock}.

\subsection*{Análisis: el patrón clásico}

La forma más simple de generar deadlock con dos mutex es que dos
tareas los adquieran en \textbf{orden inverso}: una toma $A$ y
después espera $B$; la otra toma $B$ y después espera $A$.

\subsection*{Código}

\begin{lstlisting}[style=cstyle]
lock_t lock_A;
lock_t lock_B;

/* Tarea T1: toma A, luego B */
void *thread1(void *arg) {
    lock(&lock_A);
    /* aqui puede haber una interrupcion / context switch */
    lock(&lock_B);          /* <-- espera B */

    /* seccion critica usando A y B */
    use(A, B);

    unlock(&lock_B);
    unlock(&lock_A);
    return NULL;
}

/* Tarea T2: toma B, luego A */
void *thread2(void *arg) {
    lock(&lock_B);
    /* aqui puede haber una interrupcion / context switch */
    lock(&lock_A);          /* <-- espera A */

    /* seccion critica usando A y B */
    use(A, B);

    unlock(&lock_A);
    unlock(&lock_B);
    return NULL;
}
\end{lstlisting}

\subsection*{Traza que produce el deadlock}

\begin{center}
\small
\begin{tabular}{@{}cll@{}}
\toprule
$t$  & $T_1$                       & $T_2$ \\
\midrule
$t_1$ & \texttt{lock(\&lock\_A)} $\to$ obtiene $A$ & --- \\
$t_2$ & (context switch)            & \\
$t_3$ & ---                         & \texttt{lock(\&lock\_B)} $\to$ obtiene $B$ \\
$t_4$ & ---                         & \texttt{lock(\&lock\_A)} $\to$ \textbf{bloqueado} \\
$t_5$ & (context switch)            & \\
$t_6$ & \texttt{lock(\&lock\_B)} $\to$ \textbf{bloqueado} & \\
\bottomrule
\end{tabular}
\end{center}

\textbf{Resultado}: $T_1$ tiene $A$ y espera $B$. $T_2$ tiene $B$ y
espera $A$. Ninguno puede progresar; ambos están dormidos esperando
un lock que el otro retiene. \textbf{Deadlock.}

\subsection*{Grafo de espera}

\begin{center}
\begin{tikzpicture}[
    every node/.style={font=\ttfamily\small},
    proc/.style={circle,draw,minimum size=0.95cm,font=\bfseries\small},
    res/.style={rectangle,draw,minimum size=0.95cm,font=\bfseries\small},
    every edge/.style={draw,-{Stealth[length=2.5mm]},thick}
]
\node[proc] (T1) at (0, 1.5) {$T_1$};
\node[proc] (T2) at (4, 1.5) {$T_2$};
\node[res]  (LA) at (-1.5, -1) {$A$};
\node[res]  (LB) at (5.5, -1)  {$B$};

\draw[->,blue!60!black,thick] (LA) -- node[left,font=\scriptsize]{tiene} (T1);
\draw[->,blue!60!black,thick] (LB) -- node[right,font=\scriptsize]{tiene} (T2);
\draw[->,red!70!black,thick,dashed] (T1) -- node[above,font=\scriptsize]{espera} (LB);
\draw[->,red!70!black,thick,dashed] (T2) -- node[above,font=\scriptsize]{espera} (LA);
\end{tikzpicture}
\end{center}

Se cumple el ciclo $T_1 \to B \to T_2 \to A \to T_1$. Las cuatro
condiciones de Coffman se satisfacen:

\begin{enumerate}[leftmargin=*,nosep]
    \item \textbf{Exclusión mutua}: cada lock es de uso exclusivo.
    \item \textbf{Hold and wait}: cada tarea tiene un lock y espera
        otro.
    \item \textbf{No preemption}: el SO no puede quitarle un lock a
        un thread.
    \item \textbf{Espera circular}: el ciclo descripto.
\end{enumerate}

\begin{respuesta}
El ejemplo presentado produce deadlock cuando $T_1$ adquiere
\texttt{lock\_A} y $T_2$ adquiere \texttt{lock\_B} antes de que
cualquiera intente tomar el segundo lock.

\textbf{Prevención estándar}: imponer un \emph{orden global} de
adquisición. Si ambas tareas tomaran siempre \texttt{lock\_A} antes
que \texttt{lock\_B} (e.g., orden por ID o por dirección de memoria),
el deadlock es imposible: la primera tarea en obtener $A$ retiene la
oportunidad de tomar también $B$ sin que nadie esté pidiendo el otro
en orden inverso.
\end{respuesta}

% =====================================================================
\section{Ejercicio 6 --- V/F sobre paginación (1 pto)}
% =====================================================================

\subsection*{6.a --- ``El paginado permite definir un espacio lógico a cada proceso''}

\begin{respuesta}
\textbf{V}. Cada proceso tiene su propia tabla de páginas. La MMU,
al hacer el cambio de contexto, carga el registro de tabla de
páginas (\texttt{CR3} en x86, \texttt{satp} en RISC-V) con la
dirección de la tabla del proceso entrante. A partir de ese momento,
las direcciones virtuales que emite la CPU se traducen según esa
tabla. Así, dos procesos pueden usar la misma dirección virtual
\texttt{0x1000} sin colisionar: cada uno mapea a un frame físico
distinto.
\end{respuesta}

\subsection*{6.b --- ``Dos procesos que referencian la misma dirección virtual están compartiendo memoria''}

\begin{respuesta}
\textbf{F}. Por defecto, dos procesos con sus tablas de páginas
\emph{independientes} mapean la \emph{misma dirección virtual} a
\emph{frames físicos distintos}. Que dos procesos lean
\texttt{0x1000} no significa nada en común: pueden estar
referenciando memoria completamente distinta.

\textbf{La memoria compartida real} ocurre cuando dos tablas de
páginas (de procesos distintos) tienen entradas que apuntan al
\emph{mismo frame físico}. Esto se logra explícitamente con
\texttt{shmget}/\texttt{shmat}, \texttt{mmap(MAP\_SHARED)}, librerías
compartidas (\texttt{.so}), etc. Y nada obliga a que las direcciones
virtuales coincidan entre procesos: típicamente \emph{no} coinciden.
\end{respuesta}

\subsection*{6.c --- ``Page fault cuando el bit V está apagado''}

\begin{respuesta}
\textbf{V}. El bit \texttt{V} (valid) en la PTE indica si la entrada
es legítima. Si \texttt{V = 0}, el hardware genera un \emph{page
fault} en cuanto la MMU intenta traducir esa dirección. Es el
mecanismo básico para detectar:

\begin{itemize}[leftmargin=*,nosep]
    \item Accesos a direcciones no asignadas (typical NULL deref).
    \item Demand paging (la página existe lógicamente pero todavía
        no se cargó del disco).
    \item Copy-on-write (la página existe pero como solo-lectura
        marcada para CoW).
    \item Crecimiento dinámico del stack.
\end{itemize}
\end{respuesta}

\subsection*{6.d --- ``Tabla de páginas implementa \texttt{v2p: virt $\to$ phys}''}

\begin{respuesta}
\textbf{V}. La función conceptual de la tabla de páginas es
exactamente esa: dado un \texttt{vpn} (los bits altos de la
dirección virtual), devolver el \texttt{ppn} correspondiente y
concatenar con el offset para formar la dirección física. Es una
función determinista mientras el contenido de la tabla no cambie.
\end{respuesta}

\subsection*{6.e --- ``El SO debería hacer panic ante un page fault de usuario''}

\begin{respuesta}
\textbf{F}. Un page fault \emph{no es un error fatal}. Es el
mecanismo principal por el cual el SO implementa varias optimizaciones
y servicios:

\begin{itemize}[leftmargin=*,nosep]
    \item \textbf{Demand paging}: solo se cargan en RAM las páginas
        efectivamente accedidas; las demás generan page fault que
        el kernel atiende cargándolas del disco.
    \item \textbf{Copy-on-write}: el fork() inicial no copia
        nada; las páginas se marcan solo-lectura. Al primer write,
        un page fault dispara la copia real.
    \item \textbf{Crecimiento dinámico del stack/heap}: cuando un
        proceso accede a una página recién asignada, el kernel la
        aloja en el momento.
    \item \textbf{Mapeos de archivos} (\texttt{mmap}): el contenido
        del archivo se trae a memoria solo cuando se referencia.
\end{itemize}

El SO debe \textbf{atender} el page fault. Si la dirección es
\emph{realmente inválida} (no pertenece a ningún VMA del proceso),
recién entonces envía \texttt{SIGSEGV} al proceso, que normalmente
termina. \textbf{Nunca debe hacer \texttt{panic}}: un panic significa
falla del kernel, no de un proceso de usuario.
\end{respuesta}

% =====================================================================
\section{Ejercicio 7 --- Paginación en arquitectura de 32 bits (2 ptos)}
% =====================================================================

\textbf{Enunciado.} En una arquitectura de 32 bits, con páginas de
$4$\,MB, los procesos usan un espacio de direcciones
\texttt{[0, 0x8000000)}. Ayuda: $4$\,MB $=$ \texttt{0x400000}.

\subsection*{Análisis de parámetros}

\begin{itemize}[leftmargin=*,nosep]
    \item Tamaño de página: $4$\,MB $= 2^{22}$ bytes.
    \item Espacio de direcciones del proceso:
        \texttt{[0, 0x8000000)} $= 2^{27}$ bytes $= 128$\,MB.
    \item Espacio direccionable de la arquitectura: $2^{32}$ bytes $=
        4$\,GB (mucho mayor que lo que usa cada proceso).
\end{itemize}

\subsection*{7.a --- Cantidad de páginas}

\textbf{Páginas físicas totales en la arquitectura}:

\quad $2^{32} \div 2^{22} = 2^{10} = 1024$ páginas físicas

\textbf{Páginas virtuales que usa el proceso}:

\quad \texttt{0x8000000} $\div$ \texttt{0x400000} $= 2^{27} \div 2^{22} = 2^5 = 32$ páginas

\begin{respuesta}
\textbf{(7.a)} En total, la arquitectura permite hasta
\textbf{1024 páginas} físicas (si toda la RAM estuviera presente).
El proceso, restringido a \texttt{[0, 0x8000000)}, usa hasta
\textbf{32 páginas virtuales}.
\end{respuesta}

\subsection*{7.b --- Interpretación de la dirección virtual}

\begin{itemize}[leftmargin=*,nosep]
    \item Total: $32$ bits.
    \item Offset (dentro de la página de $4$\,MB): $22$ bits bajos.
    \item Page index (vpn): $32 - 22 = 10$ bits altos.
\end{itemize}

\begin{center}
\begin{tikzpicture}[every node/.style={font=\ttfamily\small},
    cell/.style={draw,minimum height=0.7cm,align=center}]
\node[cell,minimum width=2.6cm,fill=c1!20] (vpn) at (0,0) {vpn};
\node[cell,minimum width=5.6cm,fill=c2!20,right=0pt of vpn] (off) {offset};
\node[font=\scriptsize] at ($(vpn.north)+(0,0.25)$) {10 bits};
\node[font=\scriptsize] at ($(off.north)+(0,0.25)$) {22 bits};
\node[font=\scriptsize] at ($(vpn.north west)+(0.1,-0.05)$) {31};
\node[font=\scriptsize] at ($(vpn.north east)+(-0.15,-0.05)$) {22};
\node[font=\scriptsize] at ($(off.north west)+(0.1,-0.05)$) {21};
\node[font=\scriptsize] at ($(off.north east)+(-0.05,-0.05)$) {0};
\end{tikzpicture}
\end{center}

\textbf{Cómo procesa la MMU una dirección virtual $va$:}

\begin{enumerate}[leftmargin=*,nosep]
    \item Extrae los $10$ bits altos $\to$ \texttt{vpn}.
    \item Accede a la entrada \texttt{vpn} de la tabla de páginas.
    \item Si \texttt{V = 0} o se violan los permisos $\to$ page fault.
    \item Si \texttt{V = 1}, concatena el \texttt{ppn} de la entrada
        con los $22$ bits del offset $\to$ dirección física.
\end{enumerate}

\subsection*{7.c --- Tabla de páginas del proceso}

El proceso usa: $2$ páginas de código, $1$ de datos globales,
$1$ libre (no mapeada), $2$ de stack, en ese orden a partir de la
dirección $0$.

\begin{itemize}[leftmargin=*,nosep]
    \item \texttt{vpn 0}: código (página 1).
    \item \texttt{vpn 1}: código (página 2).
    \item \texttt{vpn 2}: datos globales.
    \item \texttt{vpn 3}: \textbf{libre (no mapeada)}.
    \item \texttt{vpn 4}: stack (página 1).
    \item \texttt{vpn 5}: stack (página 2).
    \item \texttt{vpn 6} en adelante: no asignadas.
\end{itemize}

Asignamos los \texttt{ppn} arbitrariamente (los elige el SO):

\begin{center}
\small
\begin{tabular}{@{}cccccp{4.5cm}@{}}
\toprule
vpn        & ppn          & V & R/W/X    & Rango virtual                        & Área \\
\midrule
$0$        & \texttt{0x10} & 1 & R + X    & \texttt{[0x0000000, 0x0400000)}      & código (página 1) \\
$1$        & \texttt{0x11} & 1 & R + X    & \texttt{[0x0400000, 0x0800000)}      & código (página 2) \\
$2$        & \texttt{0x20} & 1 & R + W    & \texttt{[0x0800000, 0x0C00000)}      & datos globales \\
$\boldsymbol{3}$ & ---     & \textbf{0} & --- & \texttt{[0x0C00000, 0x1000000)} & \textbf{libre (no mapeada)} \\
$4$        & \texttt{0x30} & 1 & R + W    & \texttt{[0x1000000, 0x1400000)}      & stack (página 1) \\
$5$        & \texttt{0x31} & 1 & R + W    & \texttt{[0x1400000, 0x1800000)}      & stack (página 2) \\
$6$--$31$  & ---           & 0 & ---      & \texttt{[0x1800000, 0x8000000)}      & no asignadas \\
\bottomrule
\end{tabular}
\end{center}

\subsection*{7.d --- ¿La dirección \texttt{0xc40000} produce page fault?}

\textbf{Descomposición} (10 bits vpn, 22 bits offset):

\quad \texttt{0xc40000} $=$ \texttt{0x00c40000} (con 32 bits)

\quad \texttt{0xc40000} $\div$ \texttt{0x400000} $= 3$ (con resto)

\quad \texttt{0xc40000} $-$ $3 \times$ \texttt{0x400000} $=$ \texttt{0xc40000} $-$ \texttt{0xc00000} $=$ \texttt{0x40000}

Por tanto: \texttt{vpn = 3}, \texttt{offset = 0x40000}.

\textbf{Consulta a la tabla}: la entrada \texttt{vpn = 3} es la página
``libre, no mapeada'' con \texttt{V = 0}.

\begin{respuesta}
\textbf{(7.d)} \textbf{Sí, produce page fault.} La dirección
\texttt{0xc40000} cae en la \texttt{vpn = 3}, que en la tabla del
proceso tiene \texttt{V = 0} (es la página intencionalmente libre
entre datos y stack). El hardware levanta page fault y, como el
proceso no espera tener esa dirección mapeada, el kernel le envía
\texttt{SIGSEGV}.
\end{respuesta}

\subsection*{7.e --- ¿A qué área pertenece \texttt{va = 0x1400000}?}

\textbf{Descomposición}:

\quad \texttt{0x1400000} $\div$ \texttt{0x400000} $= 5$ (sin resto)

\quad $\Rightarrow$ \texttt{vpn = 5}, \texttt{offset = 0}.

\textbf{Consulta a la tabla}: la entrada \texttt{vpn = 5} corresponde
a la segunda página del stack.

\begin{respuesta}
\textbf{(7.e)} \texttt{0x1400000} pertenece al área de \textbf{stack}
(segunda página). Más concretamente, es el primer byte de la segunda
página de stack: \texttt{vpn = 5}, \texttt{offset = 0}.

\textbf{Nota}: el enunciado dice ``código, texto o stack''. Como
``código'' y ``texto'' son el mismo concepto (la sección
\texttt{.text} del ejecutable), suponemos que la opción intermedia
debía decir ``datos''. La respuesta correcta es claramente
\textbf{stack}.
\end{respuesta}

% =====================================================================
\section*{Resumen del parcial}
\addcontentsline{toc}{section}{Resumen del parcial}
% =====================================================================

\begin{center}
\small
\begin{tabular}{@{}clp{8.5cm}@{}}
\toprule
\# & Tema & Respuesta sintética \\
\midrule
1 & Microkernel IPC      & 2 traps (\texttt{send} en ambos lados) + 2 ctx switches; monolítico: 1 trap, 0 ctx \\
2 & Comandos y syscalls  & a:V, b:F (buffer del pipe), c: \texttt{sort F | uniq > final}, d:F (\texttt{exec} no crea) \\
3 & locks no atómicos    & Sin atomicidad: traza interrupciones entre test y set. En SMP con IRQs off: traza simultánea en dos CPUs \\
4 & MLFQ V/F             & a:F (no CPU-bound), b:V (sí I/O-bound), c:F (no es throughput), d:V (sí latencia interactiva) \\
5 & Deadlock             & T1 toma A,B; T2 toma B,A en orden inverso. 4 condiciones de Coffman. Fix: orden global \\
6 & Paginación V/F       & a:V, b:F (tablas distintas), c:V, d:V, e:F (no panic, atender el fault) \\
7 & Paginación numérica  & a: 1024 físicas / 32 virtuales; b: 10 + 22 bits; c: vpns 0,1,2,4,5 mapeadas y 3 libre; d: page fault; e: stack \\
\bottomrule
\end{tabular}
\end{center}

\end{document}
